% Generated draft: requires mathematical and source review.
\documentclass[11pt]{article}
\usepackage[margin=25mm]{geometry}
\usepackage{fontspec}
\setmainfont{Latin Modern Roman}
\usepackage{amsmath,amssymb,mathtools,unicode-math}
\setmathfont{Latin Modern Math}
\usepackage{enumitem,needspace,xurl,hyperref}
\setlength{\parindent}{0pt}
\setlength{\parskip}{5pt}
\setlength{\emergencystretch}{3em}
\newcommand{\N}{\mathbb{N}}
\newcommand{\Z}{\mathbb{Z}}
\newcommand{\Q}{\mathbb{Q}}
\newcommand{\R}{\mathbb{R}}
\newcommand{\C}{\mathbb{C}}
\newcommand{\F}{\mathbb{F}}
\newcommand{\A}{\mathbb{A}}
\newcommand{\PP}{\mathbb{P}}
\newcommand{\E}{\mathbb{E}}
\newcommand{\eps}{\varepsilon}
\newcommand{\dd}{\,\mathrm{d}}
\begin{document}
\section*{UnsolvedMath ID: 30004104}
\textbf{Source problem code:} catalogue entry
\subsection{Definitions and mathematical statement}
The supplied record does not define the parameter \(q\), the intended meaning of a quadrilateral, or the relation that a \(q\)-analogue must bear to that object. In particular, it specifies no ambient geometry, incidence axioms, admissible values of \(q\), or equivalence relation for candidate objects.

The source poses the following definition-seeking question:
\[
\text{What should be regarded as a }q\text{-analogue of a quadrilateral?}
\]
A precise formulation would require specifying the underlying notion of quadrilateral and the properties or correspondence that characterize its proposed \(q\)-analogue. These requirements are unresolved in the supplied evidence, so no particular construction, domain, or quantified existence or classification assertion is imposed here. The background mentions related \(q\)-Fano-type incidence structures but supplies no mathematical condition connecting them to the requested definition.
\subsection{Short English statement}
The problem asks for an appropriate definition of a \(q\)-analogue of a quadrilateral. The supplied record leaves both the underlying quadrilateral and the criteria for being an analogue unspecified, so a precise existence or classification problem cannot yet be extracted.
\subsection{Sources}
\begin{itemize}
\item Supplied local UnsolvedMath registry record, permanent ID \(30004104\), title \(q\)-Analogues of Quadrilaterals; source problem code \texttt{catalogue entry}, in the Oberwolfach Reports Open Problems collection. No individual proposer is identified; the supplied proposed year is 2019.
\item Supplied report citation: \emph{Contemporary Coding Theory (2020)}, Oberwolfach Reports, DOI 10.4171/owr/2019/13, pp.\ 773--840. \url{https://doi.org/10.4171/owr/2019/13}. The background attributes the question to workshop open-problem discussions documented in Oberwolfach Reports. The report is supplied as the source citation; its contents and the question's precise location have not been independently checked here.
\item Supplied literature reference: \emph{q-analogues of designs}. \url{https://arxiv.org/abs/1305.6126}.
\item Supplied literature reference: \emph{q-Fano plane existence}. \url{https://arxiv.org/abs/1504.06688}.
\item Supplied literature assessment, checked \textbf{2026-08-22}: the record labels the problem open and reports that the appropriate notion of a quadrilateral and related \(q\)-Fano-type incidence structures remain unresolved design questions, with no canonical general construction or classification found in that review. This is supplied dated curation, described as literature triage derived from the supplied OWR curation, and is not an independently verified assessment. Its time-sensitive status and cited claims require independent verification.
\item The supplied statement metadata labels the statement \texttt{corrected\_verified} and says that mathematical content was preserved while notation, typography, duplicate clauses, or omitted definitions were repaired for self-contained publication. This records the supplied editorial qualification and does not establish a precise mathematical formulation or independent verification here.
\end{itemize}
\end{document}
